\documentclass[a4paper]{article}

% Packages
\usepackage[utf8]{inputenc}
\usepackage{geometry}
\geometry{left=1.5cm, right=1.5cm, top=2.54cm, bottom=2.54cm}
\usepackage{graphicx, setspace, amsmath, amssymb, titlesec, fancyhdr, parskip, indentfirst, etoolbox, caption, cite, xcolor}
\usepackage{listings}
\usepackage{url}
\usepackage{hyperref}
\usepackage{multicol}
\usepackage[brazil]{babel}
\usepackage{array, longtable, colortbl, float}
\usepackage{tikz}
\usetikzlibrary{calc}

% Floats maiores no topo da página, para a EAP e as tabelas não ficarem em página própria
\renewcommand{\topfraction}{0.9}
\renewcommand{\textfraction}{0.1}
\renewcommand{\floatpagefraction}{0.8}
\makeatletter
\setlength{\@fptop}{0pt}
\setlength{\@fpsep}{16pt}
\setlength{\@fpbot}{0pt plus 1fil}
\makeatother

% Cores da EAP, usadas também no cabeçalho das tabelas
\definecolor{eapA}{HTML}{0B4F7C}
\definecolor{eapB}{HTML}{1D82C4}
\definecolor{eapC}{HTML}{E8F2FB}

\addto\captionsbrazil{
  \renewcommand{\refname}{Referências}
}

% Title Formatting
\titleformat{\section}{\centering\large\scshape}{\thesection}{1em}{}
\titleformat{\subsection}{\normalsize\bfseries}{\thesubsection.}{1em}{}

\setstretch{1.0}
\setlength{\parskip}{6pt}

\titlespacing{\section}{0pt}{14pt}{6pt}
\titlespacing{\subsection}{0pt}{6pt}{6pt}
\titlespacing{\subsubsection}{0pt}{6pt}{6pt}

% Document Title
\title{
    \textbf{Grupo 7 - Exercício 1 (versão final)}
}

\date{}

\captionsetup{labelfont={small,sc}, textfont={small,sc}}

% Section Numbering
\renewcommand{\thesection}{\Roman{section}.}
\renewcommand{\thesubsection}{\textit{\Alph{subsection}.}}
\renewcommand{\thesubsubsection}{\textit{\arabic{subsubsection}.}}
\renewcommand{\thetable}{\Roman{table}}
\renewcommand{\thefigure}{\Roman{figure}}

% Make titles italic as well
\titleformat{\subsection}
    {\normalfont\large\itshape}
    {\thesubsection}
    {1em}
    {}

\titleformat{\subsubsection}
    {\normalfont\itshape}
    {\thesubsubsection}
    {1em}
    {}

\setcounter{page}{1}

% Fancy Header Configuration
\pagestyle{fancy}
\fancyhf{}

% First Page Header
\fancypagestyle{firstpage}{
    \fancyhead[C]{
        \centering
        {\fontsize{14pt}{10pt}\selectfont
        \textbf{Universidade de São Paulo}\\[0.5em]
        \textbf{ACH2027 - Gestão de Projetos de Tecnologia da Informação}}\\[0.2em]
        {\fontsize{8pt}{10pt}\selectfont
        \textbf{Prof. Doutor:} Edmir Parada Vasques Prado \hspace{1cm}
        \textbf{Ano:} 2º Semestre/2026}
    }
}

% Default Header for Other Pages
\fancyhead[C]{
    \textbf{Escola de Artes, Ciências e Humanidades - Universidade de São Paulo}
}

\begin{document}

\maketitle
\vspace{-1.5cm}
\thispagestyle{firstpage}

% =========================================================
% INTEGRANTES
% =========================================================

\begin{multicols}{2}
    \centering

    \textbf{Gustavo G. França do Nascimento}\\
    15637551\\
    T94
    
    \vspace{0.5cm}
    
    \textbf{Kauã Lima de Souza}\\
    15674702\\
    T94
    
    \vspace{0.5cm}

    \textbf{Kevin Rodrigues Nunes}\\
    15676030\\
    T94
    
    \vspace{0.5cm}
    
    \textbf{Victor Yodono}\\
    13829040\\
    T94

\end{multicols}

\vspace{-0.2cm} % Ajuste de espaçamento vertical

% =========================================================
% CONTEÚDO
% =========================================================

\singlespacing
\setlength{\parskip}{6pt}
\setlength{\parindent}{0.5cm}

\section{Declaração de Escopo do Projeto}

De acordo com o \textbf{Guia PMBOK} (Capítulo 5 - Gerenciamento do Escopo do Projeto), a Declaração do Escopo do Projeto descreve detalhadamente as entregas do projeto e o trabalho necessário para criá-las. Ela fornece um entendimento comum do escopo entre as partes interessadas e contém as restrições, premissas e exclusões.

Abaixo está a Declaração de Escopo para o \textbf{Sistema de Controle de Acesso da EACH}.

\subsection{Descrição do Escopo do Produto}

Desenvolvimento e implantação de um sistema (software e equipamentos associados) para gerenciar e controlar a entrada e saída de pessoas e veículos nas dependências da EACH. O sistema automatizará o acesso da comunidade acadêmica (docentes, funcionários, alunos) e terceirizados regulares, além de fornecer um módulo para o cadastro, registro e liberação de visitantes. A solução contará com controle de hierarquia de acesso via senha, módulos especiais para gestão de grandes eventos e um painel para geração de relatórios gerenciais e de controle de frequência.

\subsection{Critérios de Aceitação}

Para que o projeto seja considerado concluído e aceito pela Diretoria da unidade, as seguintes condições devem ser atendidas:

\begin{itemize}
    \item \textbf{Funcionais:} O sistema deve liberar o acesso automaticamente para usuários cadastrados (via Nº USP ou CPF) e reter visitantes até que o cadastro seja efetuado.
    \item \textbf{Segurança e Regras:} O acesso ao sistema de gestão deve exigir autenticação por senha e respeitar níveis hierárquicos pré-definidos (ex: administrador, operador de portaria).
    \item \textbf{Desempenho (Relatórios):} O sistema deve ser capaz de emitir, sem falhas, relatórios exatos de frequência de entrada de pessoas (filtrado por Nº USP e CPF) e veículos (filtrado por categoria).
    \item \textbf{Operacionais:} As rotinas especiais (como o pré-cadastramento em lote para eventos de grande porte) devem funcionar adequadamente, sem sobrecarregar a portaria.
    \item \textbf{Cronograma e Custos:} O projeto deve ser totalmente entregue e implantado no prazo máximo de 1 (um) ano, respeitando os limites orçamentários estabelecidos para equipamentos e recursos humanos.
\end{itemize}

\subsection{Entregas do Projeto (Deliverables)}

\begin{itemize}
    \item Documento de Requisitos e Arquitetura do Sistema.
    \item \textbf{Módulo de Controle de Acesso:} Funcionalidades de leitura e liberação de catracas/cancelas para pessoas e veículos.
    \item \textbf{Módulo de Gestão de Usuários e Visitantes:} Interface para cadastro de perfis (docentes, alunos, funcionários, terceirizados e visitantes) e definição de hierarquias de acesso.
    \item \textbf{Módulo de Rotinas Especiais:} Área do sistema dedicada ao pré-cadastro rápido/em lote para grandes eventos.
    \item \textbf{Módulo de Relatórios:} Painel gerador de relatórios de controle e frequência (por Nº USP, CPF e Veículos).
    \item Equipamentos e materiais de controle de acesso (catracas, cancelas, leitores, etc.) comprados, instalados e configurados.
    \item Treinamento para a equipe de operação (portaria/segurança) e administração do sistema.
\end{itemize}

\subsection{Exclusões do Projeto (O que está fora do escopo)}

\begin{itemize}
    \item Controle de acesso interno de portas de salas de aula, laboratórios e gabinetes (o foco é a entrada principal/perímetro da Universidade).
    \item Desenvolvimento de módulos de cobrança (tarifação de estacionamento).
    \item Manutenção e suporte técnico contínuo do sistema após a implantação final e encerramento do projeto.
    \item Sistemas de câmeras de segurança (CFTV) não integrados diretamente à catraca/cancela.
\end{itemize}

\subsection{Restrições (Constraints)}

Limitações impostas ao projeto, relativas a tempo, custo, recursos e regulamentos:

\begin{itemize}
    \item \textbf{Prazo:} A implantação total do sistema não pode ultrapassar 1 (um) ano. O desenvolvimento deve ser finalizado em agosto de 2026.
    \item \textbf{Orçamento (Equipamentos):} O gasto com aquisição de equipamentos e materiais está limitado à faixa de R\$ 300.000,00 a R\$ 600.000,00.
    \item \textbf{Orçamento (Bolsas):} O custo com alunos está restrito ao fornecimento de 5 a 8 bolsas de R\$ 2.000,00 mensais cada, pelo período exato de 6 meses.
    \item \textbf{Recursos Humanos:} O desenvolvimento do sistema é restrito à participação de alunos do último ano de SI (onde a atividade contará como estágio obrigatório) \textbf{OU} à contratação de uma Start-Up da incubadora da EACH.
    \item A equipe exige a presença de um Líder de Projeto subordinado à Diretoria e um Analista de Sistemas Sênior (graduado em SI).
\end{itemize}

\subsection{Premissas (Assumptions)}

Fatores considerados verdadeiros, reais ou certos para fins de planejamento:

\begin{itemize}
    \item A administração da Universidade (USP/EACH) fornecerá acesso seguro e integrado à base de dados atual de docentes, funcionários e alunos (Números USP, CPFs e dados de veículos) para a carga inicial do sistema.
    \item As instâncias acadêmicas competentes aprovarão a participação dos alunos no projeto para validação dos créditos de estágio de graduação a tempo do início em agosto de 2026.
    \item Caso a opção escolhida seja o uso de alunos bolsistas, haverá alunos de SI suficientes, capacitados e interessados no último ano letivo para preencher as 5 a 8 vagas.
    \item A infraestrutura física e de rede da EACH nos pontos de entrada (portarias) suportará a instalação dos novos equipamentos sem necessidade de obras civis complexas e fora do orçamento previsto.
\end{itemize}

\section{Estrutura Analítica do Projeto (EAP)}

A Estrutura Analítica do Projeto (EAP) apresenta a decomposição hierárquica do trabalho necessário para atingir os objetivos do projeto \cite{pmi2017}. A versão final (Figura~\ref{fig:eap}) tem 6 grupos de entregas no segundo nível e 22 pacotes de trabalho no terceiro. Cada pacote é uma entrega, escrita como substantivo; as atividades que produzem cada pacote ficam no cronograma (Exercício 2).

\begin{figure}[htbp]
    \centering
    \begin{tikzpicture}[
            raiz/.style={fill=eapA, text=white, font=\sffamily\small\bfseries, inner sep=4pt, rounded corners=1pt},
            grupo/.style={fill=eapB, text=white, font=\sffamily\scriptsize\bfseries, text width=2.45cm,
                          minimum height=9mm, align=left, inner sep=3pt, rounded corners=1pt},
            pacote/.style={draw=eapA, fill=eapC, font=\sffamily\scriptsize, text width=2.2cm,
                           minimum height=7.5mm, align=left, inner sep=3pt, rounded corners=1pt},
            every node/.append style={execute at begin node={\hyphenpenalty=10000\exhyphenpenalty=10000}},
            liga/.style={draw=eapA, semithick}]
        \node[raiz, anchor=south] (raiz) at (8.71cm, 1.1cm) {1. Sistema de Controle de Acesso à EACH};
        \draw[liga] (raiz.south) -- (8.71cm, 0.45cm);
        \draw[liga] (1.33cm, 0.45cm) -- (16.08cm, 0.45cm);

        \node[grupo, anchor=north west] (g1) at (0.00cm, 0) {1.1 Gerenciamento do projeto};
        \draw[liga] (1.33cm, 0.45cm) -- (g1.north);
        \node[pacote, anchor=north west] (p11) at ($(g1.south west)+(0.25cm,-2mm)$) {1.1.1 Linha de base do escopo};
        \draw[liga] ($(g1.south west)+(0.12cm,0)$) |- (p11.west);
        \node[pacote, anchor=north west] (p12) at ($(p11.south west)+(0,-1.6mm)$) {1.1.2 Cronograma};
        \draw[liga] ($(g1.south west)+(0.12cm,0)$) |- (p12.west);
        \node[pacote, anchor=north west] (p13) at ($(p12.south west)+(0,-1.6mm)$) {1.1.3 Orçamento};
        \draw[liga] ($(g1.south west)+(0.12cm,0)$) |- (p13.west);
        \node[pacote, anchor=north west] (p14) at ($(p13.south west)+(0,-1.6mm)$) {1.1.4 Equipe do projeto};
        \draw[liga] ($(g1.south west)+(0.12cm,0)$) |- (p14.west);
        \node[pacote, anchor=north west] (p15) at ($(p14.south west)+(0,-1.6mm)$) {1.1.5 Relatórios de acompanhamento};
        \draw[liga] ($(g1.south west)+(0.12cm,0)$) |- (p15.west);
        \node[pacote, anchor=north west] (p16) at ($(p15.south west)+(0,-1.6mm)$) {1.1.6 Encerramento};
        \draw[liga] ($(g1.south west)+(0.12cm,0)$) |- (p16.west);

        \node[grupo, anchor=north west] (g2) at (2.95cm, 0) {1.2 Requisitos};
        \draw[liga] (4.28cm, 0.45cm) -- (g2.north);
        \node[pacote, anchor=north west] (p21) at ($(g2.south west)+(0.25cm,-2mm)$) {1.2.1 Documento de requisitos};
        \draw[liga] ($(g2.south west)+(0.12cm,0)$) |- (p21.west);
        \node[pacote, anchor=north west] (p22) at ($(p21.south west)+(0,-1.6mm)$) {1.2.2 Documento de arquitetura};
        \draw[liga] ($(g2.south west)+(0.12cm,0)$) |- (p22.west);
        \node[pacote, anchor=north west] (p23) at ($(p22.south west)+(0,-1.6mm)$) {1.2.3 Protótipos de tela};
        \draw[liga] ($(g2.south west)+(0.12cm,0)$) |- (p23.west);

        \node[grupo, anchor=north west] (g3) at (5.90cm, 0) {1.3 Infraestrutura};
        \draw[liga] (7.23cm, 0.45cm) -- (g3.north);
        \node[pacote, anchor=north west] (p31) at ($(g3.south west)+(0.25cm,-2mm)$) {1.3.1 Equipamentos adquiridos};
        \draw[liga] ($(g3.south west)+(0.12cm,0)$) |- (p31.west);
        \node[pacote, anchor=north west] (p32) at ($(p31.south west)+(0,-1.6mm)$) {1.3.2 Equipamentos instalados};
        \draw[liga] ($(g3.south west)+(0.12cm,0)$) |- (p32.west);
        \node[pacote, anchor=north west] (p33) at ($(p32.south west)+(0,-1.6mm)$) {1.3.3 Rede das portarias};
        \draw[liga] ($(g3.south west)+(0.12cm,0)$) |- (p33.west);

        \node[grupo, anchor=north west] (g4) at (8.85cm, 0) {1.4 Desenvolvimento};
        \draw[liga] (10.18cm, 0.45cm) -- (g4.north);
        \node[pacote, anchor=north west] (p41) at ($(g4.south west)+(0.25cm,-2mm)$) {1.4.1 Módulo de Controle de Acesso};
        \draw[liga] ($(g4.south west)+(0.12cm,0)$) |- (p41.west);
        \node[pacote, anchor=north west] (p42) at ($(p41.south west)+(0,-1.6mm)$) {1.4.2 Módulo de Usuários e Visitantes};
        \draw[liga] ($(g4.south west)+(0.12cm,0)$) |- (p42.west);
        \node[pacote, anchor=north west] (p43) at ($(p42.south west)+(0,-1.6mm)$) {1.4.3 Módulo de Rotinas Especiais};
        \draw[liga] ($(g4.south west)+(0.12cm,0)$) |- (p43.west);
        \node[pacote, anchor=north west] (p44) at ($(p43.south west)+(0,-1.6mm)$) {1.4.4 Módulo de Relatórios};
        \draw[liga] ($(g4.south west)+(0.12cm,0)$) |- (p44.west);

        \node[grupo, anchor=north west] (g5) at (11.80cm, 0) {1.5 Testes};
        \draw[liga] (13.13cm, 0.45cm) -- (g5.north);
        \node[pacote, anchor=north west] (p51) at ($(g5.south west)+(0.25cm,-2mm)$) {1.5.1 Testes unitários};
        \draw[liga] ($(g5.south west)+(0.12cm,0)$) |- (p51.west);
        \node[pacote, anchor=north west] (p52) at ($(p51.south west)+(0,-1.6mm)$) {1.5.2 Testes de integração};
        \draw[liga] ($(g5.south west)+(0.12cm,0)$) |- (p52.west);
        \node[pacote, anchor=north west] (p53) at ($(p52.south west)+(0,-1.6mm)$) {1.5.3 Homologação};
        \draw[liga] ($(g5.south west)+(0.12cm,0)$) |- (p53.west);

        \node[grupo, anchor=north west] (g6) at (14.75cm, 0) {1.6 Implantação};
        \draw[liga] (16.08cm, 0.45cm) -- (g6.north);
        \node[pacote, anchor=north west] (p61) at ($(g6.south west)+(0.25cm,-2mm)$) {1.6.1 Carga inicial de dados};
        \draw[liga] ($(g6.south west)+(0.12cm,0)$) |- (p61.west);
        \node[pacote, anchor=north west] (p62) at ($(p61.south west)+(0,-1.6mm)$) {1.6.2 Treinamento};
        \draw[liga] ($(g6.south west)+(0.12cm,0)$) |- (p62.west);
        \node[pacote, anchor=north west] (p63) at ($(p62.south west)+(0,-1.6mm)$) {1.6.3 Sistema em produção};
        \draw[liga] ($(g6.south west)+(0.12cm,0)$) |- (p63.west);
    \end{tikzpicture}
    \caption{Estrutura Analítica do Projeto (EAP), versão final}
    \label{fig:eap}
\end{figure}

\subsection{Mudanças em relação ao esboço}

A revisão aplicou três regras do PMBOK \cite{pmi2017}: a regra dos 100\% (a EAP cobre todo o trabalho do projeto e nada além dele), a orientação a entregas (substantivos, não verbos) e um único responsável por pacote de trabalho. A Tabela~\ref{tab:mudancas} resume o que mudou.

\begin{table}[htbp]
    \centering
    \small
    \renewcommand{\arraystretch}{1.3}
    \begin{tabular}{|>{\raggedright\arraybackslash}p{4cm}|>{\raggedright\arraybackslash}p{5cm}|>{\raggedright\arraybackslash}p{7.1cm}|}
        \hline
        \rowcolor{eapC}
        \textbf{Esboço} & \textbf{Versão final} & \textbf{Motivo} \\
        \hline
        1.1.1 Escopo & 1.1.1 Linha de base do escopo & O elemento passa a nomear a entrega, e não o assunto. \\
        \hline
        1.1.3 Reuniões & 1.1.5 Relatórios de acompanhamento & Reunião é trabalho contínuo, e a entrega é o relatório. O grupo já tinha decidido tirar as reuniões (Registro 01 do Exercício 2). \\
        \hline
        (não havia) & 1.1.3 Orçamento, 1.1.4 Equipe do projeto e 1.1.6 Encerramento & Regra dos 100\%: respeitar o orçamento, contratar a equipe e encerrar o projeto também são trabalho do projeto. \\
        \hline
        1.2.1 Especificação, 1.2.2 Arquitetura, 1.2.3 Telas & Documento de requisitos, Documento de arquitetura, Protótipos de tela & Nomes de entregas. \\
        \hline
        1.3 Infra/Hardware: Aquisição, Instalação, Rede & 1.3 Infraestrutura: Equipamentos adquiridos, Equipamentos instalados, Rede das portarias & Nomes de entregas. \\
        \hline
        1.4.1 Integração, 1.4.2 Usuários, 1.4.3 Eventos, 1.4.4 Relatórios & Módulos com os nomes da declaração de escopo & Cada entrega da declaração aparece com o mesmo nome na EAP. \\
        \hline
        1.5.1 Unitários, 1.5.2 Integração & Testes unitários, Testes de integração & Clareza. \\
        \hline
        1.6.1 Carga de Dados, 1.6.3 Go-Live & Carga inicial de dados, Sistema em produção & Nome em português e como entrega. \\
        \hline
    \end{tabular}
    \caption{Mudanças da EAP em relação ao esboço.}
    \label{tab:mudancas}
\end{table}

% O dicionário começa em página nova, depois da EAP e da tabela de mudanças
\clearpage
\section{Dicionário da EAP}

O dicionário (Tabela~\ref{tab:dicionario}) descreve cada pacote de trabalho da EAP \cite{pmi2017}. Cada pacote tem um único responsável, entre os dois papéis que a declaração de escopo define: o Líder de Projeto e o Analista de Sistemas Sênior. As durações de cada pacote estão no cronograma (Exercício 2), e os custos serão detalhados na gestão de custos.

{\small
\renewcommand{\arraystretch}{1.3}
\begin{longtable}{|>{\raggedright\arraybackslash}p{2.9cm}|>{\raggedright\arraybackslash}p{5.3cm}|>{\raggedright\arraybackslash}p{5.3cm}|>{\raggedright\arraybackslash}p{2.2cm}|}
    \hline
    \rowcolor{eapC}
    \textbf{Pacote de trabalho} & \textbf{Descrição do trabalho} & \textbf{Critério de aceitação} & \textbf{Responsável} \\
    \hline
    \endhead
    \rowcolor{eapC!45}
    \multicolumn{4}{|l|}{\textbf{1.1 Gerenciamento do projeto}} \\*
    \hline
    1.1.1 Linha de base do escopo & Declaração do escopo, EAP e dicionário da EAP. & Aprovada pela Diretoria. Mudanças depois disso passam pelo controle de mudanças. & Líder de Projeto \\
    \hline
    1.1.2 Cronograma & Atividades, durações, precedências e rede do projeto. & Aprovado pela Diretoria, com a implantação total em até 1 ano. & Líder de Projeto \\
    \hline
    1.1.3 Orçamento & Custos de equipamentos e de bolsas, ou do contrato com a start-up. & Equipamentos entre R\$ 300.000,00 e R\$ 600.000,00; de 5 a 8 bolsas de R\$ 2.000,00 por 6 meses. & Líder de Projeto \\
    \hline
    1.1.4 Equipe do projeto & Seleção de 5 a 8 alunos do último ano de SI, com o estágio aprovado, ou contratação de uma start-up da incubadora da EACH. & Equipe formalizada antes do início do desenvolvimento. & Líder de Projeto \\
    \hline
    1.1.5 Relatórios de acompanhamento & Relatórios de andamento, custos e riscos para a Diretoria. & Entregues na periodicidade combinada com a Diretoria. & Líder de Projeto \\
    \hline
    1.1.6 Encerramento & Termo de encerramento e registro das lições aprendidas. & Termo assinado pela Diretoria depois da entrada em produção. O suporte posterior fica fora do escopo. & Líder de Projeto \\
    \hline
    \rowcolor{eapC!45}
    \multicolumn{4}{|l|}{\textbf{1.2 Requisitos}} \\*
    \hline
    1.2.1 Documento de requisitos & Regras de negócio e casos de uso dos quatro módulos, levantados com a Diretoria e a portaria. & Aprovado pelo Líder de Projeto e pela Diretoria. & Analista Sênior \\
    \hline
    1.2.2 Documento de arquitetura & Tecnologias, integração com catracas, cancelas e leitores, e acesso à base de dados da USP. & Acesso à base da USP confirmado com a área responsável por ela. & Analista Sênior \\
    \hline
    1.2.3 Protótipos de tela & Telas de cadastro, de portaria e de relatórios. & Validados com operadores da portaria antes do desenvolvimento. & Analista Sênior \\
    \hline
    \rowcolor{eapC!45}
    \multicolumn{4}{|l|}{\textbf{1.3 Infraestrutura}} \\*
    \hline
    1.3.1 Equipamentos adquiridos & Especificação, cotação e compra de catracas, cancelas e leitores. & Itens recebidos conforme a especificação e dentro da faixa de orçamento. & Líder de Projeto \\
    \hline
    1.3.2 Equipamentos instalados & Preparação dos pontos de entrada e fixação dos equipamentos, sem obras civis complexas. & Equipamentos fixados e ligados em todos os pontos de entrada. & Analista Sênior \\
    \hline
    1.3.3 Rede das portarias & Cabeamento e configuração da rede entre os equipamentos e o servidor. & Todos os equipamentos se comunicam com o servidor. & Analista Sênior \\
    \hline
    \rowcolor{eapC!45}
    \multicolumn{4}{|l|}{\textbf{1.4 Desenvolvimento}} \\*
    \hline
    1.4.1 Módulo de Controle de Acesso & Comunicação com os equipamentos e rotina de liberação para pessoas e veículos. & Libera automaticamente quem está cadastrado (Nº USP ou CPF) e retém o visitante até o cadastro. & Analista Sênior \\
    \hline
    1.4.2 Módulo de Usuários e Visitantes & Cadastro de docentes, funcionários, alunos, terceirizados e visitantes, e níveis de acesso ao sistema. & Acesso ao sistema de gestão só com senha, respeitando os níveis (ex.: administrador, operador de portaria). & Analista Sênior \\
    \hline
    1.4.3 Módulo de Rotinas Especiais & Pré-cadastro em lote para eventos de grande porte e emissão das credenciais do evento. & Evento de grande porte cadastrado em lote, sem sobrecarregar a portaria. & Analista Sênior \\
    \hline
    1.4.4 Módulo de Relatórios & Consultas e painel gerencial de frequência. & Relatórios exatos de frequência de pessoas (por Nº USP e CPF) e de veículos (por categoria). & Analista Sênior \\
    \hline
    \rowcolor{eapC!45}
    \multicolumn{4}{|l|}{\textbf{1.5 Testes}} \\*
    \hline
    1.5.1 Testes unitários & Scripts de teste dos quatro módulos e correção dos defeitos encontrados. & Todos os testes passando antes dos testes de integração. & Analista Sênior \\
    \hline
    1.5.2 Testes de integração & Passagens simuladas nos equipamentos instalados, inclusive com falha de rede. & Liberação e retenção corretas em todos os pontos de entrada. & Analista Sênior \\
    \hline
    1.5.3 Homologação & Validação do sistema pela Diretoria, com os critérios de aceitação da declaração de escopo. & Termo de aceite assinado pela Diretoria. & Líder de Projeto \\
    \hline
    \rowcolor{eapC!45}
    \multicolumn{4}{|l|}{\textbf{1.6 Implantação}} \\*
    \hline
    1.6.1 Carga inicial de dados & Extração e importação dos dados de docentes, funcionários, alunos e veículos da base da USP. & Dados importados e conferidos por amostragem. & Analista Sênior \\
    \hline
    1.6.2 Treinamento & Manuais e aulas práticas para a portaria, a segurança e os administradores. & Equipe de operação treinada e manuais entregues. & Analista Sênior \\
    \hline
    1.6.3 Sistema em produção & Ativação do sistema em todas as portarias. & Sistema operando em todas as portarias, com os dados da carga inicial. & Líder de Projeto \\
    \hline
    \caption{Dicionário da EAP.}
    \label{tab:dicionario}
\end{longtable}
}

\section{Registro de Uso de Inteligência Artificial}

Em conformidade com a Política de Uso de Ferramentas de Inteligência Artificial Generativa da disciplina ACH2027, declaramos explicitamente a utilização de IA como suporte ao aprendizado, ideação e estruturação deste documento.

\begin{table}[H]
    \centering
    \renewcommand{\arraystretch}{1.5}
    \small
    \begin{tabular}{|p{3cm}|p{6cm}|p{6cm}|}
        \hline
        \rowcolor{eapC}
        \textbf{Registro de Uso de IA} & \textbf{Registro 01: Declaração de Escopo e LaTeX} & \textbf{Registro 02: Diagramação da EAP (Mermaid)} \\
        \hline
        \textbf{Problema / Objetivo} & Elaborar a Declaração de Escopo (PMBOK) a partir do caso semestral da atividade proposta e formatar o conteúdo no template LaTeX do grupo. & Criar o diagrama visual da Estrutura Analítica do Projeto (EAP) e gerar o código exportável para o Overleaf. \\
        \hline
        \textbf{Ferramenta de IA} & Gemini 3.1 Pro (Setembro/2026). & Gemini 3.1 Pro (Setembro/2026). \\
        \hline
        \textbf{Comandos} & "Com base no projeto a seguir, elabore a declaração de escopo... do PMBOK" e "Formate isso no seguinte latex...". & "faça um esboço da estrutura analítica... precisa ser um diagrama" e "me de um mermaid... ficou muito largo horizontalmente". \\
        \hline
        \textbf{Revisão dos resultados} & A primeira versão não abordava detalhes técnicos de software. Perguntamos se era necessário e recebemos um tetalhamento de questões como Open Source e LGPD. Decidimos não incluir, já que a declaração do escopo é uma visão alto nível de gestão. & O formato padrão (Top-Down) do Mermaid extrapolava as margens da página. Ajustamos as instruções para gerar formatos otimizados (Left-Right) adequados ao papel A4 no Overleaf. \\
        \hline
        \textbf{Aplicação no trabalho} & Seção de Declaração de Escopo do Projeto. & Seção da Estrutura Analítica do Projeto (EAP). \\
        \hline
    \end{tabular}
    \caption{Registro de Uso de IA Generativa.}
    \label{tab:uso_ia}
\end{table}

\begin{table}[H]
    \centering
    \renewcommand{\arraystretch}{1.5}
    \small
    \begin{tabular}{|p{3cm}|p{12.4cm}|}
        \hline
        \rowcolor{eapC}
        \textbf{Registro de Uso de IA} & \textbf{Registro 03: EAP final e dicionário da EAP} \\
        \hline
        \textbf{Problema / Objetivo} & Revisar o esboço da EAP com as boas práticas do PMBOK e produzir a versão final da EAP e o dicionário da EAP. \\
        \hline
        \textbf{Ferramenta de IA} & Claude Opus 5.5, via Claude Code (Setembro/2026). \\
        \hline
        \textbf{Comandos} & "pode fazer a EAP final e o dicionário, mas crie um arquivo novo a partir da entrega inicial, para compararmos depois" e "pode refinar o layout dos docs". \\
        \hline
        \textbf{Revisão dos resultados} & \textcolor{red}{[PREENCHER: revisão do grupo]} \\
        \hline
        \textbf{Aplicação no trabalho} & Seções II (EAP e mudanças em relação ao esboço) e III (dicionário da EAP), e o acabamento visual do documento. \\
        \hline
    \end{tabular}
    \caption{Registro de Uso de IA Generativa (continuação).}
    \label{tab:uso_ia_final}
\end{table}

\begin{thebibliography}{99}

\bibitem{pmi2021}
PMI - PROJECT MANAGEMENT INSTITUTE. \textbf{Um Guia do Conhecimento em Gerenciamento de Projetos (Guia PMBOK)}. 7. ed. Pensilvânia: Project Management Institute, 2021.

\bibitem{pmi2017}
PMI - PROJECT MANAGEMENT INSTITUTE. \textbf{Um Guia do Conhecimento em Gerenciamento de Projetos (Guia PMBOK)}. 6. ed. Newtown Square: Project Management Institute, 2017.

\end{thebibliography}

\end{document}